% Compile with: pdflatex -output-directory=/tmp/llm-report llm-response.tex
\documentclass[11pt,letterpaper]{article}
\usepackage[T1]{fontenc}
\usepackage{lmodern}
\usepackage[margin=1in]{geometry}
\usepackage{microtype}
\usepackage{xcolor}
\usepackage{enumitem}
\usepackage{amsmath}
\usepackage[most]{tcolorbox}
\usepackage{hyperref}

\definecolor{ink}{HTML}{243244}
\definecolor{accent}{HTML}{245D73}
\definecolor{paperblue}{HTML}{F2F6F8}
\definecolor{linegray}{HTML}{D6DEE4}
\hypersetup{colorlinks=true,urlcolor=accent,linkcolor=accent,
  pdftitle={Research report: prompt and LLM response}}
\color{ink}
\setlength{\parindent}{0pt}
\setlength{\parskip}{0.65em}
\linespread{1.12}
\setlist{leftmargin=*,itemsep=0.3em,topsep=0.3em}
\pagestyle{plain}
\tcbset{message/.style={enhanced,breakable,sharp corners,
  boxrule=0.5pt,colframe=linegray,colback=white,
  coltext=ink,coltitle=ink,fonttitle=\sffamily\bfseries,
  colbacktitle=paperblue,left=12pt,right=12pt,top=10pt,bottom=10pt,
  before skip=12pt,after skip=12pt}}
\newtcolorbox{promptmessage}{message,title={PROMPT},colback=paperblue}
\newtcolorbox{responsemessage}{message,title={LLM RESPONSE}}

\begin{document}

{\small\sffamily\bfseries\color{accent} START RESEARCH REPORT}\par
{\LARGE\bfseries Research report title}\par
{\small Topic: [topic]\quad Model: [model/version]\quad Date: [date]}\par
\medskip
\hrule height 0.5pt

% Replace the example below with your prompt. Use LaTeX, not raw Markdown.
\begin{promptmessage}
Explain [research question] for [intended reader]. State the main answer
first, explain the reasoning, and identify the evidence and limitations.

\textbf{Context.} [Background, definitions, and relevant constraints.]

\textbf{Requested outline.}
\begin{enumerate}
  \item Main answer and key findings.
  \item Analysis, evidence, and sources.
  \item Limitations and next steps.
\end{enumerate}
\end{promptmessage}

% Repeat the two message environments for additional turns.
% Both boxes can continue across pages for long prompts and responses.
\begin{responsemessage}
\textbf{Main answer.} [Write the central answer in two or three sentences.
Make the most useful conclusion easy to find.]

\subsection*{Key findings}
\begin{itemize}
  \item \textbf{Finding one.} [State the finding and its significance.]
  \item \textbf{Finding two.} [Add supporting detail or a comparison.]
\end{itemize}

\subsection*{Analysis and evidence}
[Develop one idea per paragraph. Explain why the evidence supports the
answer. Define unfamiliar terms before using them.]

% Place verified citations beside claims; use \href{URL}{Source title} for links.

\subsection*{Limitations}
[Separate observations, interpretation, and uncertainty. Identify what
cannot be concluded from the available evidence.]

\subsection*{Next steps}
\begin{enumerate}
  \item {[First concrete action or follow-up question.]}
  \item {[Second action and the uncertainty it resolves.]}
\end{enumerate}
\end{responsemessage}

\end{document}
